\section{Conclusion}
\label{sec:conclusion}

We presented the \textbf{Failure-Trace Allocator (FTA)}, a deterministic,
gold-free answer-selection policy for matched-budget LLM inference.
On the primary Aggregate-720 evaluation (three disjoint seeds, zero overlap),
FTA achieves 80.69\% (581/720) versus L1 77.64\%, S1 77.08\%, and TALE
75.14\%.
Source-stratified 95\% bootstrap confidence intervals are strictly positive
for all four external baseline comparisons, with the lower bound versus the
best external at $+1.11$~pp.
On the Final-300 evaluation (seed~71, $n=300$), FTA reaches 86.67\% (260/300),
compared to 83.00\% for L1, 82.00\% for S1, and 78.33\% for TALE.

The policy consists of two rule-based gates applied in strict priority order.
The Low-Depth Guard (LDG) detects structurally weak frontier search branches
using logged failure-trace metadata and substitutes the external majority
answer.
The External-Consensus Fallback (ECO), a low-frequency high-precision
safeguard in this evaluation, identifies cases where all three external
baselines unanimously disagree with the frontier's selection, overriding
to the consensus answer.
Both gates are gold-free: no correctness labels, example identifiers, or
artifact paths are accessed at runtime, and no additional inference calls are
issued after candidate generation.
The rules, trigger conditions, and exact field names are fully specified in
Table~\ref{tab:policy_rules} and Table~\ref{alg:fix24}, with reproducibility
details documented in the appendix.

As supporting preliminary evidence, a multi-provider regime-learning experiment (D9) using
the same failure-trace signal achieves a cross-validated accuracy of
$0.5018 \pm 0.0252$ against a frontier baseline of $0.3436$, with zero
false-override events across all tested providers and thresholds.
These preliminary results are consistent with the failure-trace signal having
provider-independent utility, though a full multi-provider evaluation at
matched scale remains for future work.

Failure analysis of the 40 residual errors on Final-300 shows that 24 (60\%)
are \emph{all-methods-wrong} cases in which no evaluated method produced the
correct answer; the dominant subtype is multi-step composition failures
(10/24) where the error is embedded in intermediate reasoning rather than in
the final selection step.
Further accuracy gains in this setting are therefore likely to require
advances in candidate generation rather than selection policy alone.

Future directions include learned allocation policies that replace the
deterministic gates with calibrated classifiers, evaluation on broader
datasets and provider families, stronger candidate-pool construction
strategies that reduce the pool-miss failure class, and extension to
multi-turn or open-ended generation settings where failure-trace signals may
take different structural forms.
